\section{本章小结}
\label{sec:pgrad-summary}

章首的问题是：样本平均$\frac{1}{n}\sum_{i}X_{i}^{2}$里看不到参数，对它求导得到零，可$\mathbb{E}[X^{2}]$明明随$\theta$变化。答案是参数通过样本的来源进入了期望，求梯度时必须把它找出来并计入。本章的做法是先把参数的位置看清，再把期望写成某个固定测度上的积分，参数全部移进被积函数，然后用定理~\ref{thm:pgrad-exchange}的局部Lipschitz控制检查求导能否移进积分。逐点可微不够，图~\ref{fig:pgrad-moving-spike}的移动尖峰就是反例。

三种表示对应三种固定测度的选法，各自利用不同的计算条件，也各有失效的情形。

\begin{BookInlineTable}
  \centering
  \caption{三种梯度表示的比较。}
  \label{tab:pgrad-methods}
  \begin{tabular}{@{}llll@{}}
    \toprule
    & 直接求导 & 重参数化 & 得分函数 \\
    \midrule
    固定的测度 & 抽样分布$P$ & 噪声分布$Q$ & 参考测度$\nu$ \\
    需要计算 & $\nabla_{\bm{\theta}}f_{\bm{\theta}}$ & 生成映射及$\nabla_{x}f_{\bm{\theta}}$ & 得分或$\bm{a}_{\bm{\theta}}$ \\
    对$f$的要求 & 对参数可控 & 沿路径可控 & 可求值即可 \\
    典型失效 & 控制不可积 & 离散选择、跳跃 & 支撑随参数移动 \\
    \bottomrule
  \end{tabular}
\end{BookInlineTable}

直接求导适用于抽样分布固定的情形，梯度就是逐样本导数的平均。重参数化把分布的变化改写成样本点的移动，路径导数由链式法则的两条路组成；函数沿路径跳跃时它看不到跳跃的贡献。得分函数估计对概率权重求导，不需要$f$对$x$可微，离散变量与不连续函数都能处理；但它要求支撑不随参数移动，否则会漏掉边界项。归一化常数未知时，得分是未归一化对数导数减去它的期望，梯度化为函数值与对数导数的协方差，只需目标样本；用独立样本估计时，以样本平均为中心的协方差要除以$n-1$。

无论用哪种表示，都要分开检查两件事：梯度样本的可积性保证估计无偏和一致，二阶矩才保证误差按$n^{-1/2}$缩小。不同表示的方差可以相差很多，没有哪一种在所有问题上更好；得分估计可以借助得分期望为零加入基线，基线若用到当前样本，须像留一基线那样保证与该样本的得分独立。样本来自近似生成器或有限步链时，估计误差分成与样本数无关的采样偏差和有限样本波动；对有限步程序求导得到的是程序期望的导数，它与目标导数是否相同要另行证明。

本章只回答怎样把期望的导数变成可估计的积分，以及这个估计的偏差与波动。用这些梯度逐步更新参数，是随机优化的问题。下一章在共同的精度要求下，比较各种采样、方差控制与梯度表示的计算预算，讨论怎样为具体任务选择方法。
